\tzpanchor{1155}
\tzpentryheader{1155}{RESERVED SIGNAL CONTINUITY GAP ID1155}

\begin{tzpsnapshot}
\tzpsnaprow{TZPID}{\tzpcode{ID1155}}
\tzpsnaprow{UUID}{\tzpcode{74b961dc-bd49-5bad-a3ce-2bfe4b5cdbb7}}
\tzpsnaprow{Type}{Transfer Statement}
\tzpsnaprow{Scale}{transfer}
\tzpsnaprow{LOC Classification}{Working routing: QA641 manifolds and topology; QC174.17 mathematical physics}
\tzpsnaprow{arXiv Classification}{Primary: math-ph; Cross-list: gr-qc, math.DG}
\end{tzpsnapshot}

\tzpsection{Canonical Statement}
This entry develops reserved signal continuity gap id1155 inside the signal arc of the directory. It treats the named object as a TRAWIN-admissible carrier that can be rendered, normalized, and linked forward into the surrounding registry pages.
\[
c30a236,registry_id_0524,[ID 0524],Delta function in the bell curve (2),,equation_like_line,"\boxed{\,\Psi(R,\mathbf r,t)=\mathcal N\,F(R)\
\]
\[
abla\times \mathbf{E}=-\partial\\_t \mathbf{B}
\]

\noindent
\begin{minipage}[t]{0.48\linewidth}
\tzpsection{Wolfram Form}
\begingroup
\small
\[
\begin{aligned}
\mathfrak{W}_{\mathrm{TZP}}\!\left(\mathcal{K}_{1155}\right)
&\longmapsto
\left\langle \Omega^{\mathbb{T}},\; \Psi_{\mathrm{T}},\; \rho_{\mathrm{vac}} \right\rangle
\end{aligned}
\]
\[
\mathcal{K}_{1155}
:=
\operatorname{HoldForm}\!\left(\mathcal{T}_{1155}\right)
\]
\[
\begin{aligned}
\operatorname{Admissible}_{\mathrm{TZP}}\!\left(\mathcal{K}_{1155}\right)
&\Longleftrightarrow
\mathfrak{W}_{\mathrm{TZP}}\!\left(\mathcal{K}_{1155}\right)\not\equiv\varnothing
\end{aligned}
\]
\textbf{Wolfram register:} canonical symbol lift\\
\textbf{Title lift:} RESERVED SIGNAL CONTINUITY GAP ID1155\\
\textbf{Symbol key:} \(\Omega^{\mathbb{T}}\): Trawinistic boundary curvature\\ \(\Psi_{\mathrm{T}}\): Trawinistic wavefunction\\ \(\rho_{\mathrm{vac}}\): vacuum energy density
\par
\endgroup
\end{minipage}\hfill
\begin{minipage}[t]{0.48\linewidth}
\tzpsection{TRAWIN Normalization}
\[
\tzpnormalizewrap{c30a236,registry_id_0524,[ID 0524],Delta function in the bell curve (2),,equation_like_line,"\boxed{\,\Psi(R,\mathbf r,t)=\mathcal N\,F(R)\,\,abla\times \mathbf{E}=-\partial\\_t \mathbf{B},\,Delta r = \lambda/(2 \text{NA}) \approx 180}
\]

\small
\textbf{T}: Variation Law is localized within the signal arc of the directory.\\
\textbf{R}: Support Relation preserves the operative relation that reserved signal continuity gap id1155 requires.\\
\textbf{A}: Closure Relation fixes the admissible closure condition for the entry.\\
\textbf{W}: RESERVED SIGNAL CONTINUITY GAP ID1155 is rendered as a reusable operator-level statement.\\
\textbf{I}: The informational content of reserved signal continuity gap id1155 is retained rather than flattened.\\
\textbf{N}: RESERVED SIGNAL CONTINUITY GAP ID1155 closes as a distinct normalized TRAWIN object.
\normalsize
\end{minipage}

\tzpsection{Derivable Consequences}
The entry now reads through actual sourced equations, so its local arc is carried by explicit mathematical relations rather than a uniform quaternary certificate record shell.
\clearpage

\tzpentryheader{1155}{Oxford Dictionary of the Queens, NY English}

\begingroup\small
\tzpsection{Definition}
RESERVED SIGNAL CONTINUITY GAP ID1155 is a registry definition in Signal with secondary pressure from Engineering and Implementation Arcs.

% BEGIN TZP CONTEXTUAL MEANING
\tzpsection{Contextual Meaning}
\begingroup\footnotesize\setlength{\emergencystretch}{3em}\sloppy
\textbf{Formal statement:} This entry develops reserved signal continuity gap id1155 inside the signal arc of the directory. It treats the named object as a TRAWIN-admissible carrier that can be rendered, normalized, and linked forward into the surrounding registry pages. \textbf{Contextual meaning:} The entry reads RESERVED SIGNAL CONTINUITY GAP ID1155 as a controlled technical headword whose equation layer, proof route, and registry role are interpreted together. \textbf{Derivable consequences:} The entry now reads through actual sourced equations, so its local arc is carried by explicit mathematical relations rather than a uniform quaternary certificate record shell. \textbf{Lemma support:} Variation Law; Support Relation; Closure Relation.\par
\endgroup
% END TZP CONTEXTUAL MEANING

\tzpsection{Technical Interpretation}
Technically, RESERVED SIGNAL CONTINUITY GAP ID1155 is read through the local equation chain and the TRAWIN normalization recorded on the entry page.

\tzpsection{Equation Grounding}
Equation anchors: partial t E; ablatimes E =-partial t B; and Delta r = lambda/(2 NA ) approximately 180. Proof route: show that the stated equations form a coherent progression for the entry and remain compatible under the TRAWIN ordering. Formalization route: prove that the final closure relation follows from the preceding entry-specific equations.

\tzpsection{Interpretive Role}
The entry now reads through actual sourced equations, so its local arc is carried by explicit mathematical relations rather than a uniform quaternary certificate record shell.
\endgroup

\tzpsection{TZP Thesaurus}
\begin{enumerate}[leftmargin=1.6em,itemsep=6pt,topsep=4pt]
\item \textbf{Reserved Signal Continuity Gap Id1155}: entry-specific alternate wording for indexing, related literature, and local formal context.
\item \textbf{Reserved Signal Continuity Formal Analogue}: entry-specific alternate wording for indexing, related literature, and local formal context.
\item \textbf{Id1155 Equivalence}: entry-specific alternate wording for indexing, related literature, and local formal context.
\end{enumerate}

\tzpsection{Encyclopedia Mathematica}
\begin{samepage}
\noindent\begin{minipage}[t]{0.45\linewidth}
\vspace{0pt}
\raggedright
\scriptsize\textbf{Image reading.} The rendered manifold at right is the per-ID light plate for RESERVED SIGNAL CONTINUITY GAP ID1155. It should be read as the visual companion to the entry's equation layer, not as a repeated template diagram.\par
\end{minipage}\hfill
\begin{minipage}[t]{0.53\linewidth}
\vspace*{-0.10in}
\raggedright
\tzpencyclopediaimage{1155}
\end{minipage}
\end{samepage}

\clearpage

\tzpentryheader{1155}{Direct Equation Progression and Proof Trajectory}

\tzpsection{Direct Equation Progression}
\begin{enumerate}[leftmargin=1.6em,itemsep=9pt,topsep=6pt]
\item \textbf{Variation Law}
\[
c30a236,registry_id_0524,[ID 0524],Delta function in the bell curve (2),,equation_like_line,"\boxed{\,\Psi(R,\mathbf r,t)=\mathcal N\,F(R)\
\]
This fixes the first explicit relation carried by the entry rather than a generic certificate record normalization.

\item \textbf{Support Relation}
\[
abla\times \mathbf{E}=-\partial\\_t \mathbf{B}
\]
This adds the next operational equation that advances the entry beyond its opening definition.

\item \textbf{Closure Relation}
\[
Delta r = \lambda/(2 \text{NA}) \approx 180
\]
This closes the progression with an actual admissibility or closure relation instead of recycling the normalization shell.
\end{enumerate}

\tzpsection{Proof}
\textbf{Proof route.} show that the stated equations form a coherent progression for the entry and remain compatible under the TRAWIN ordering.

\textbf{Formal method.} prove that the final closure relation follows from the preceding entry-specific equations.

\medskip
\tzpproofartifacts{Isabelle audit: no standalone source in corpus.}{Lean audit: no standalone source in corpus.}{Rocq audit: no standalone source in corpus.}

\tzpsection{Dependencies and Test Handle}
\begin{tzpsnapshot}
\tzpsnaprow{Dependencies}{\tzpidref{1154}, \tzpidref{1100}}
\tzpsnaprow{Node}{\tzpcode{registry_id_1155}}
\tzpsnaprow{Support Titles}{Variation Law; Support Relation; Closure Relation}
\tzpsnaprow{Test Handle}{If the cited entry equations cannot be made mutually admissible in one TRAWIN-normalized frame, the entry fails.}
\end{tzpsnapshot}

\tzpsection{Canonical Equation Links}
\tzpequationlinks
  {c30a236,registry_id_0524,[ID 0524],Delta function in the bell curve (2),,equation_like_line,"\boxed{\,\Psi(R,\mathbf r,t)=\mathcal N\,F(R)\}
  {\mathcal{J}_{\mathrm{out}} = \cdots}
  {\mathcal{A}_{\mathrm{adm}}\!\left(\mathcal{J}_{\mathrm{out}}\right) = \cdots}
